\chapter{Buying Connectivity}\label{ch:nw:buying-connectivity}

A cloud provider's \omterm{def:nw:buying-connectivity:sla}{service-level agreement} for its private connections promises 99.99\% of each month to customers who buy two
connections at two sites, and a credit of 10\% of the month's port charges if it fails to deliver. For one connection it promises 95\%, and
pays nothing unless the month has lost more than a day and a half. Trading firms buy circuits on the same kind of terms: the provider
promises a percentage, pays a fraction of a monthly fee when it misses, and the loss from the outage (the positions that could not be
hedged, the quotes that could not be pulled) is the firm's. This chapter is about what a firm buys, from whom, on which terms, and why the
second circuit is worth buying only if it does not share a trench with the first.

Chapters~9 to~14 described the buildings, the maps and the routes. Here they become purchases: \omterm{def:nw:colocation-products-and-how-they-are-sold:mmr}{cross-connects} and circuits in a
building, \omterm{def:nw:buying-connectivity:metro}{metro circuits} between buildings, long-haul fibre and radio between cities, and feeds delivered by others. The prices that are
public come from venues' fee schedules; the rest are negotiated, and the chapter says so.

\section{What is bought: circuits, wavelengths, bandwidth and feeds}

\begin{definition}[Metro circuit]\label{def:nw:buying-connectivity:metro}
A \emph{metro circuit}\index{metro circuit} is a point-to-point connection between two buildings in the same metropolitan area (two data centres,
a data centre and an office), sold as a \omterm{def:nw:long-haul-fibre:dark}{lit service} of a given speed or built on leased or owned fibre.
\end{definition}

A firm's connectivity is a stack of purchases. In each building it buys \omterm{def:nw:colocation-products-and-how-they-are-sold:mmr}{cross-connects} (chapter~9) and the venue's own ports; between the
buildings of a metropolitan area, \omterm{def:nw:buying-connectivity:metro}{metro circuits}; between cities, wavelengths or \omterm{def:nw:long-haul-fibre:dark}{dark fibre} (chapter~13) and radio bandwidth
(chapter~14); and, for data it does not collect itself, feeds delivered by an extranet or a vendor. A New Jersey firm that trades
the NYSE, Nasdaq and Cboe from one building buys at least two \omterm{def:nw:buying-connectivity:metro}{metro circuits} and the venues' ports at the other end; one that trades
Chicago futures against them adds a long-haul route and its back-up.

\begin{omfigure}
\begin{tikzpicture}[font=\footnotesize,b/.style={draw,rounded corners=2pt,minimum height=0.6cm,align=center,inner sep=3pt},>=Stealth]
\node[b,fill=omDef!15,minimum width=2.4cm] (f) at (0,0) {firm's cabinet\\ (Carteret)};
\node[b,fill=omThm!15,minimum width=2.2cm] (v1) at (0,2.0) {Nasdaq ports\\ (same building)};
\node[b,fill=omThm!15,minimum width=2.2cm] (v2) at (6.4,1.2) {NYSE group\\ (Mahwah)};
\node[b,fill=omThm!15,minimum width=2.2cm] (v3) at (6.4,-0.8) {Cboe\\ (Secaucus NY5)};
\node[b,fill=omThm!15,minimum width=2.2cm] (v4) at (11.6,1.2) {CME Globex\\ (Aurora)};
\node[b,fill=gray!10,minimum width=2.2cm] (x) at (0,-2.0) {extranet or\\ vendor feed};
\draw[->,thick] (f) -- node[right,font=\scriptsize]{cross-connect} (v1);
\draw[->,thick] (f) -- node[sloped,above,font=\scriptsize]{metro circuit} (v2);
\draw[->,thick] (f) -- node[sloped,below,font=\scriptsize]{metro circuit} (v3);
\draw[->,thick,omMeth] (v2.east) -- node[above,font=\scriptsize,omMeth,align=center]{long haul:\\ radio, fibre} (v4.west);
\draw[->,thick,gray] (f) -- node[right,font=\scriptsize]{circuit} (x);
\end{tikzpicture}

\omcaption{A New Jersey firm's connectivity as purchases, schematically: \omterm{def:nw:colocation-products-and-how-they-are-sold:mmr}{cross-connects} and ports in its own building, \omterm{def:nw:buying-connectivity:metro}{metro circuits} to the other
two buildings of the triangle (chapter~10), a long-haul route with its back-up to Chicago (chapters~13 and~14), and a circuit to a provider for the
data it does not collect itself.}\label{fig:nw:buying-connectivity:stack}
\end{omfigure}

\begin{dated}{2026-09}{Published connectivity prices: orders of magnitude}\label{dat:nw:buying-connectivity:prices}
The NYSE group's connectivity fee schedule (last updated April 2025) charges, per month: 2\,500~USD for a 1-gigabit IP-network circuit, 11\,000
and 18\,000~USD for 10- and 40-gigabit IP and NMS network connections, 600~USD for a fibre \omterm{def:nw:colocation-products-and-how-they-are-sold:mmr}{cross-connect} in the data centre, 9\,000 to
44\,000~USD for 10 to 200~megabits of wireless bandwidth between Mahwah and Secaucus (10\,000 to 45\,000 to Carteret), and 6\,000~USD for a
wireless connection carrying CME Group data; one-time charges of 500 to 15\,000~USD apply. MIAX Pearl charges 15\,000~USD a month for a
10-gigabit ultra-low-latency connection (chapter~9); Deutsche Börse 6\,000 to 8\,400~EUR a month for its 10-gigabit co-location connections
(chapter~11).
\end{dated}

\begin{table}[ht]
\centering\small
\begin{tabular}{lrrl}
\toprule
Item & Monthly & One-time & Source \\
 & (USD) & (USD) & \\
\midrule
NYSE IP network, 1 Gb circuit & 2\,500 & 2\,500 & fee schedule \\
NYSE IP and NMS networks, 10 Gb & 11\,000 & 10\,000 & fee schedule \\
NYSE IP and NMS networks, 40 Gb & 18\,000 & 10\,000 & fee schedule \\
NYSE data-centre fibre cross connect & 600 & 500 & fee schedule \\
NYSE wireless Mahwah--Secaucus, 10 Mb & 9\,000 & 10\,000 & fee schedule \\
NYSE wireless Mahwah--Secaucus, 200 Mb & 44\,000 & 10\,000 & fee schedule \\
NYSE wireless connection of CME data & 6\,000 & 5\,000 & fee schedule \\
MIAX Pearl 10 Gb ultra-low-latency & 15\,000 & 0 & fee filing \\
\bottomrule
\end{tabular}
\caption{The catalogue of \texttt{firm.slamodel}: published prices only, each with its ledger row and date. Wireless bandwidth costs about as much
per megabit a month as a 10-gigabit fibre connection costs per gigabit. Data: \texttt{nw\_buy.price\_rows()}.}\label{tab:nw:buying-connectivity:prices}
\end{table}

The table is the scale of the market. A megabit of radio between two New Jersey buildings costs about 900~USD a month at 10~megabits and
220~USD at 200; a gigabit of fibre between a cabinet and the venue's network costs about 1\,100~USD a month at 10~gigabits. Speed, not
bandwidth, is what the radio sells.

\section{Financial extranets and managed infrastructure}

\begin{definition}[Financial extranet]\label{def:nw:buying-connectivity:extranet}
A \emph{financial extranet}\index{financial extranet} is a private network, run by a provider connected to many venues and firms, over which its
customers reach venues' market data and order entry, and each other, without building their own circuits to each; venues admit such providers
under an agreement and charge them for the access they resell.
\end{definition}

Nasdaq's definition shows the arrangement from the venue's side: an extranet provider signs Nasdaq's extranet provider agreement, has its own
connection to Nasdaq, is co-located in a Nasdaq facility, and gives its own customers access to Nasdaq's services, in the building or elsewhere.
For a firm the extranet is a trade: one connection to the provider replaces many to the venues, at the price of the provider's latency and of
sharing its network with its other customers.

\begin{definition}[Managed infrastructure]\label{def:nw:buying-connectivity:managed}
\emph{Managed infrastructure}\index{managed infrastructure} is equipment and connectivity that a provider installs, owns or leases, and operates
on a firm's behalf (servers, switches, circuits, market-data handlers), usually in the provider's racks near the venues, sold as a monthly
service.
\end{definition}

\omterm{def:nw:buying-connectivity:managed}{Managed infrastructure} is the build-or-buy question of One Quant Book~16 (build against buy) at the scale of a rack: the firm that buys it gets
a presence at a venue in weeks instead of months, and hands the provider control over the latency of every piece in the path. Firms whose
edge is latency buy the building's space and build the rest; firms whose edge is elsewhere buy the rest too.

\section{Wireless-bandwidth vendors}

Radio routes are sold in three ways. Venues sell bandwidth or data on their own links: ICE between its Mahwah data centre and the other New
Jersey buildings (with Anova Financial Networks) and from Toronto; SIX between Zurich and Frankfurt; Euronext from London to Bergamo with
McKay Brothers. Network operators sell bandwidth on theirs. And data vendors sell a market-data feed delivered over radio, as Quincy Data did
from Aurora in chapter~10. What they have in common is scarcity: megabits, not gigabits; a fixed set of end points; and a fibre back-up that
the firm must also buy, or accept to do without when it rains.

\section{Contracts and service levels}

\begin{definition}[Service-level agreement, service credit]\label{def:nw:buying-connectivity:sla}
A \emph{service-level agreement}\index{service-level agreement} (SLA) is the part of a supply contract that states the performance the provider
commits to (availability, latency, repair time), how it is measured, and what the provider owes when it fails. A \emph{service credit}\index{service
credit} is what it owes: usually a percentage of the month's charges for the failed service, credited against future bills.
\end{definition}

\begin{definition}[Mean time between failures, mean time to repair]\label{def:nw:buying-connectivity:mtbf}
The \emph{mean time between failures}\index{mean time between failures} (MTBF) of a component is the average time it runs between failures; its
\emph{mean time to repair}\index{mean time to repair} (MTTR) is the average time from a failure to its restoration. Its steady-state availability is
$\mathrm{MTBF}/(\mathrm{MTBF}+\mathrm{MTTR})$.
\end{definition}

\begin{dated}{2026-09}{A published credit schedule}\label{dat:nw:buying-connectivity:sla}
AWS's \omterm{def:nw:buying-connectivity:sla}{service-level agreement} for Direct Connect (last updated July 2026) makes three commitments: 99.99\% monthly uptime for multi-site
redundant deployments, 99.9\% for multi-site non-redundant ones, and 95.0\% for a single connection. Missing them earns a credit of 10\%, 25\% or
100\% of the affected connections' port charges, at thresholds of (99.99, 99.0, 95.0), (99.9, 99.0, 95.0) and (95.0, 92.5, 90.0)\% respectively.
Credits are the sole remedy, and single connections are not recommended for production.
\end{dated}

\begin{definition}[Route diversity]\label{def:nw:buying-connectivity:diversity}
\emph{Route diversity}\index{route diversity} is the property of two or more circuits between the same end points that they share no physical
element whose failure would stop both: no duct, cable, building entry, room, piece of equipment or power supply.
\end{definition}

\begin{proposition}[Two circuits]\label{prop:nw:buying-connectivity:two}
Let each of two circuits be unavailable a fraction $u$ of the time, independently, and let common-mode events that stop both occur a fraction
$c$ of the time. The pair is unavailable a fraction $u^2 + c - u^2 c \approx u^2 + c$. When $c$ is of the order of $u$, the second circuit
gains almost nothing; only diversity, $c \to 0$, delivers $u^2$.
\end{proposition}

\begin{proof}
The pair is down when both are down independently (probability $u^2$) or a common-mode event is under way ($c$); inclusion--exclusion gives the
union.
\end{proof}

\omcode{code/firm/slamodel/firm_slamodel.py}{78}{84}{A design's steady-state availability: every circuit down at once, or a common-mode event.}\label{lst:nw:buying-connectivity:avail}

\begin{omfigure}
\begin{tikzpicture}[font=\footnotesize,b/.style={draw,rounded corners=2pt,minimum height=0.55cm,align=center,inner sep=2pt},>=Stealth]
\foreach \y/\lab in {2.4/{one circuit},0.9/{two circuits in one duct},-0.8/{two diverse circuits}}{
  \node[b,fill=omDef!15,minimum width=1.6cm] at (0,\y) {firm};
  \node[b,fill=omThm!15,minimum width=1.6cm] at (9,\y) {venue};
  \node[anchor=west,font=\scriptsize] at (-1,\y+0.55) {\lab};}
\draw[thick] (0.8,2.4) -- (8.2,2.4);
\draw[gray,line width=4pt,opacity=0.3] (2.2,0.9) -- (6.8,0.9);
\draw[thick] (0.8,1.0) -- (8.2,1.0); \draw[thick] (0.8,0.8) -- (8.2,0.8);
\node[gray,font=\scriptsize] at (4.5,0.5) {one duct: one digger cuts both};
\draw[thick] (0.8,-0.7) -- (2,-0.2) -- (7,-0.2) -- (8.2,-0.7);
\draw[thick] (0.8,-0.9) -- (2,-1.4) -- (7,-1.4) -- (8.2,-0.9);
\node[font=\scriptsize] at (4.5,0.05) {route A}; \node[font=\scriptsize] at (4.5,-1.65) {route B: another street, entry, provider};
\end{tikzpicture}

\omcaption{Three designs, schematically: one circuit; two circuits whose fibres share a duct, so that one event (a cut, a fire) stops both; and two
diverse circuits that share nothing between the firm's and the venue's rooms.}\label{fig:nw:buying-connectivity:designs}
\end{omfigure}

The model gives each circuit two failures a year with a mean repair of four hours (MTBF \qty{4383}{\hour}, MTTR \qty{4}{\hour}), and the shared
duct half a common-mode event a year with a mean repair of twelve hours (a cut fibre is spliced, not swapped): assumptions, labelled as such.
One circuit is then down 0.091\% of the year, 480 minutes. Two in one duct are down 360 minutes: the duct's six hours a year, almost all of it,
since two independent failures coincide for less than half a minute a year. Two diverse circuits are down 0.4 minutes a year.

\begin{table}[ht]
\centering\small
\begin{tabular}{lrrrrr}
\toprule
Design & Availability & Expected & 95th pct. & Annual cost & Loss at 2\,000 \\
 & (\%) & (min/yr) & (min/yr) & (kUSD) & USD/min (kUSD) \\
\midrule
One circuit & 99.9088 & 479.6 & 1\,541 & 139.2 & 959 \\
Two circuits in one duct & 99.9315 & 360.4 & 1\,992 & 278.4 & 721 \\
Two diverse circuits & 99.99992 & 0.4 & 0 & 320.2 & 0.9 \\
\bottomrule
\end{tabular}
\caption{Simulation of three designs of 10-gigabit NYSE connections (published prices; the failure rates, the duct's common-mode events, a 15\%
premium for diversity and the trading loss per minute are assumptions): steady-state availability and expected minutes down, the 95th percentile
of a simulated year (1\,000 years), annual cost with one \omterm{def:nw:colocation-products-and-how-they-are-sold:mmr}{cross-connect} per circuit, and the expected trading loss. Data:
\texttt{nw\_buy.results()}.}\label{tab:nw:buying-connectivity:designs}
\end{table}

\begin{omfigure}
\begin{tikzpicture}
\begin{semilogyaxis}[width=10.5cm,height=5.8cm,xmin=100,xmax=440,ymin=0.1,ymax=2000,xlabel={annual cost (thousand USD)},
  ylabel={expected minutes down a year (log)},tick label style={font=\footnotesize},label style={font=\small},grid=major,grid style={gray!25},
  table/col sep=comma,nodes near coords,point meta=explicit symbolic,every node near coord/.append style={font=\scriptsize,anchor=west}]
\addplot[only marks,mark=*,mark size=2.6pt,omThm] table[x=cost_k,y=expected_min,meta=design]{figdata/networks/15-buying-connectivity/designs.csv};
\end{semilogyaxis}
\end{tikzpicture}

\omcaption{Expected minutes down a year against annual cost for the three designs of \cref{tab:nw:buying-connectivity:designs}: the second circuit in
the same duct doubles the cost and removes a quarter of the downtime; diversity adds 15\% and removes nearly all of it. Data: \texttt{fig\_buy.py}.}\label{fig:nw:buying-connectivity:cost}
\end{omfigure}

\omcode{code/firm/slamodel/firm_slamodel.py}{118}{129}{A simulated year: each circuit's outages, the common-mode events, and the minutes during
which every circuit was down at once.}\label{lst:nw:buying-connectivity:sim}

\begin{omfigure}
\begin{tikzpicture}
\begin{semilogyaxis}[width=11cm,height=6cm,xmin=0,xmax=100,ymin=0.01,ymax=20000,xlabel={percentile of simulated years},
  ylabel={minutes down in the year (log)},tick label style={font=\footnotesize},label style={font=\small},grid=major,grid style={gray!25},
  legend columns=3,legend style={font=\footnotesize,at={(0.5,-0.25)},anchor=north,draw=none},table/col sep=comma]
\addplot[omDef,thick,no markers] table[x=pct,y=minutes]{figdata/networks/15-buying-connectivity/ecdf_single.csv};
\addplot[omMeth,thick,dashed,no markers] table[x=pct,y=minutes]{figdata/networks/15-buying-connectivity/ecdf_duct.csv};
\addplot[omThm,thick,no markers] table[x=pct,y=minutes]{figdata/networks/15-buying-connectivity/ecdf_diverse.csv};
\legend{one circuit,two circuits in one duct,two diverse circuits}
\end{semilogyaxis}
\end{tikzpicture}

\omcaption{Simulation: minutes down in a year by percentile of 1\,000 simulated years (zero drawn at 0.01). Two circuits in one duct have more
years without an outage than one circuit, and worse bad years: a duct event takes twelve hours on average. Diverse circuits almost never go
down together. Data: \texttt{fig\_buy.py}.}\label{fig:nw:buying-connectivity:ecdf}
\end{omfigure}

What do the \omterm{def:nw:buying-connectivity:sla}{service credits} pay for these outages? Applying the published schedule of \cref{dat:nw:buying-connectivity:sla} to each simulated month,
the single circuit, on the single-connection terms, earns an average of 5.50~USD a year: a four-hour outage leaves the month at 99.44\%, well above
95\%. The two circuits in one duct, on the non-redundant terms, earn about 2\,560~USD a year; the diverse pair, on the redundant terms, nothing.
Against losses of the order of hundreds of thousands of dollars in the model, credits are not insurance: they are a signal of what the provider
expects to deliver, and the firm insures itself by design.

\section{Orders of magnitude for price}

The published prices put a floor under a budget. A firm with two diverse 10-gigabit connections to the NYSE group pays about 320\,000~USD a year
for them in the model; the same firm's two ultra-low-latency connections to an options exchange cost 360\,000~USD (chapter~9); a radio route of 200
megabits between two New Jersey buildings costs more than half a million a year. What is not published (extranet services, \omterm{def:nw:buying-connectivity:metro}{metro circuits} between
private buildings, long-haul wavelengths, managed racks) is priced by quotation, and chapter~29 treats those rows as estimates with ranges, never as
facts.

\begin{method}[Buying a circuit]\label{met:nw:buying-connectivity:buy}
\begin{enumerate}
\item Say what the circuit is for: which strategies, what latency, what bandwidth, what happens when it fails.
\item Ask each provider for the path (ducts, streets, building entries, rooms), the one-way rack-to-rack latency and how it is measured, and its
  outage history.
\item Buy diversity, not duplicates: check that the second circuit shares nothing with the first, down to the building entry and the power.
\item Read the SLA for what it measures (monthly uptime, repair time), the credits and the exclusions; compare the credits with your own loss per
  minute.
\item Put every price, term and renewal date in the budget with its source (chapter~29), and re-check before renewal.
\end{enumerate}
\end{method}

\section{Tutorial: a year of outages}\label{tut:nw:buying-connectivity:tut}

\begin{tutorial}
\textbf{Goal.} Price three designs from published fees, simulate their outages, and compare credits with losses.
\textbf{End state:} \cref{fig:nw:buying-connectivity:ecdf} and \cref{tab:nw:buying-connectivity:prices,tab:nw:buying-connectivity:designs}.
\begin{tutsteps}
\item \textbf{The catalogue.} \texttt{firm\_slamodel.CATALOGUE} holds published prices with their ledger rows; \texttt{SCHEDULES} the published
  credit tiers.
\item \textbf{The designs.} \texttt{nw\_buy.designs()} builds one circuit, two in one duct and two diverse ones from \texttt{ASSUME}.
\item \textbf{Steady state and simulation.} \texttt{design\_availability} (\cref{lst:nw:buying-connectivity:avail}) and \texttt{simulate\_year}
  (\cref{lst:nw:buying-connectivity:sim}); \texttt{results()} fills \cref{tab:nw:buying-connectivity:designs}.
\item \textbf{Credits.} \texttt{monthly\_credits} cuts each simulated year into months and applies the schedule.
\end{tutsteps}
\textbf{What to change next.} Make the duct events rarer but longer; put the two diverse circuits with the same provider and add a shared
provider fault; change the loss per minute to your own strategy's figure.
\end{tutorial}

\section{Build: the SLA model}\label{bld:nw:buying-connectivity:slamodel}

\begin{build}
\bfield{Purpose} Circuits as priced and sourced catalogue rows, their availability alone and in designs, the credits a contract pays and the losses it
does not: the input of chapters~28 (resilience) and~29 (the plan).
\bfield{Interface} \texttt{firm\_slamodel}: \texttt{Tier}, \texttt{Circuit}, \texttt{Design}, \texttt{CATALOGUE}, \texttt{SCHEDULES},
\texttt{availability}, \texttt{design\_availability}, \texttt{simulate\_year}, \texttt{monthly\_down\_min}, \texttt{uptime\_pct}, \texttt{credit},
\texttt{downtime\_cost}, \texttt{expected\_down\_min}, \texttt{nines}.
\bfield{Rules} Prices come only from the catalogue, with their source and date; failure rates are stated assumptions; a design's common-mode events
are explicit, never zero by default when circuits share anything.
\bfield{Acceptance tests} \texttt{code/firm/slamodel/tests/}: availability by hand for the three designs, simulation against the steady state,
monthly bucketing, and the credit tiers.
\bfield{Stretch} Repair times that depend on the failure (a cut against a card), maintenance windows announced in advance, and a provider's own
common mode across its customers.
\end{build}

\begin{omsources}
\item AWS Direct Connect service-level agreement (July 2026).
\item Nasdaq Data News \#2017-8 (extranet providers); NYSE group connectivity fee schedule (April 2025).
\item Chapters~9 to~14 for the venues' and vendors' own descriptions of their connections.
\end{omsources}

\section{Exercises}

\begin{exercise}[$\star$]\label{exo:nw:buying-connectivity:1}
A circuit fails twice a year and takes four hours to repair each time. What is its availability, and how many minutes is it down a year?
\end{exercise}

\begin{exercise}[$\star$]\label{exo:nw:buying-connectivity:2}
From \cref{tab:nw:buying-connectivity:prices}, what does a megabit of wireless bandwidth between Mahwah and Secaucus cost a month at 10 and at 200
megabits?
\end{exercise}

\begin{exercise}[$\star$]\label{exo:nw:buying-connectivity:3}
Under the single-connection schedule of \cref{dat:nw:buying-connectivity:sla}, how long must a connection be down in a 30-day month before any credit
is due?
\end{exercise}

\begin{exercise}[$\star\star$]\label{exo:nw:buying-connectivity:4}
Using \cref{prop:nw:buying-connectivity:two}, why do two circuits in one duct gain so little over one?
\end{exercise}

\begin{exercise}[$\star\star$]\label{exo:nw:buying-connectivity:5}
List five things two ``diverse'' circuits can still share.
\end{exercise}

\begin{exercise}[$\star\star$]\label{exo:nw:buying-connectivity:6}
When is a \omterm{def:nw:buying-connectivity:extranet}{financial extranet} the right way to reach a venue, and when is it not?
\end{exercise}

\begin{exercise}[$\star\star\star$]\label{exo:nw:buying-connectivity:7}
\emph{Coding.} With \texttt{firm.slamodel}, find the common-mode rate at which two circuits in one duct are no better than one circuit.
\end{exercise}

\begin{exercise}[$\star\star\star$]\label{exo:nw:buying-connectivity:8}
\emph{Find the flaw.} ``Our provider guarantees 99.99\% and pays credits, so an outage costs us nothing.''
\end{exercise}

\section{Problem: Two Circuits Are Not Twice as Good}

\begin{problem}[{Weekend problem --- redundancy that is not}]\label{pb:nw:buying-connectivity:1}
A firm connects to the NYSE group over 10-gigabit connections. Each fails twice a year and takes four hours on average to repair. Its two
circuits run in one duct, which suffers half a common-mode event a year with twelve hours of repair; a diverse pair would cost 15\% more. The firm
loses 2\,000~USD for each minute without a connection.

\medskip\noindent\textbf{Part I --- One circuit.}
\begin{enumerate}
\item What are the circuit's availability and expected minutes down a year?
\item What does it cost a year with its \omterm{def:nw:colocation-products-and-how-they-are-sold:mmr}{cross-connect}?
\item What is the expected annual loss?
\item What credit does the single-connection schedule pay for a four-hour outage?
\end{enumerate}

\medskip\noindent\textbf{Part II --- Two circuits in one duct.}
\begin{enumerate}[resume]
\item What are the pair's availability and expected minutes down?
\item How much of that is the duct?
\item What does the pair cost, and what loss does it avoid compared with one circuit?
\item What does a bad year look like (95th percentile)?
\end{enumerate}

\medskip\noindent\textbf{Part III --- Two diverse circuits.}
\begin{enumerate}[resume]
\item What are the pair's availability and expected minutes down?
\item What does diversity cost over the duct, and what does it save?
\item What credits do the three designs earn a year on average?
\item What would make the diverse pair fail together anyway?
\end{enumerate}

\medskip\noindent\textbf{Part IV --- The verdict.}
\begin{enumerate}[resume]
\item State the \emph{named result}: the availability of two circuits sharing a duct against two diverse ones, their expected annual minutes down,
  and the \omterm{def:nw:buying-connectivity:sla}{service credit} against the trading loss.
\item Which assumption moves the answer most?
\item How would you verify a provider's claim of diversity?
\item What should the SLA measure for a trading firm, beyond monthly uptime?
\item How do the firm's systems use two circuits: active and standby, or both active?
\item Where does this analysis go in the connectivity plan of chapter~29?
\item What does the second circuit buy if it shares the duct?
\item In one sentence: what is redundancy?
\end{enumerate}
\end{problem}

\section{Interview questions}

\begin{interviewq}[$\star$ \iqroles{developer}]\label{iq:nw:buying-connectivity:1}
What is an SLA, and what does a \omterm{def:nw:buying-connectivity:sla}{service credit} compensate?
\end{interviewq}

\begin{interviewq}[$\star\star$ \iqroles{developer}]\label{iq:nw:buying-connectivity:2}
You have two circuits to the same venue from two providers. How can they still fail together?
\end{interviewq}

\begin{interviewq}[$\star\star$ \iqroles{developer, trader}]\label{iq:nw:buying-connectivity:3}
When would you use a \omterm{def:nw:buying-connectivity:extranet}{financial extranet} instead of a direct connection to a venue?
\end{interviewq}

\begin{interviewq}[$\star\star$ \iqroles{developer}]\label{iq:nw:buying-connectivity:4}
How do MTBF and MTTR combine into availability, and what does a second independent circuit do to it?
\end{interviewq}

\begin{interviewq}[$\star\star$ \iqroles{trader}]\label{iq:nw:buying-connectivity:5}
Radio bandwidth costs hundreds of times more per megabit than fibre. Why do firms buy it?
\end{interviewq}

\begin{interviewq}[$\star\star\star$ \iqroles{developer, researcher}]\label{iq:nw:buying-connectivity:6}
Design the connectivity of a new trading office to three venues in two cities for 99.999\% availability. What do you buy and what do you verify?
\end{interviewq}
